\subsection{李代数的最大幂零理想}

设 \(\mathfrak{g}\) 是李代数，\(\mathfrak{a}\) 是 \(\mathfrak{g}\) 的理想。\(\mathfrak{a}\) 幂零的充要条件是：对所有 \(x \in \mathfrak{a}\)，\(\mathrm{ad}_{\mathfrak{g}}x\) 都是幂零的；充分性显然，而必要性是因为若 \(\mathfrak{a}\) 幂零且 \(x \in \mathfrak{a}\)，则 \(\mathrm{ad}_{\mathfrak{a}}x\) 是幂零的，且 \(\mathrm{ad}_{\mathfrak{a}}x\) 将 \(\mathfrak{g}\) 映入 \(\mathfrak{a}\)，所以 \(\mathrm{ad}_{\mathfrak{g}}x\) 是幂零的。将命题 4 应用于 \(\mathfrak{g}\) 的伴随表示，得到以下结果：

命题 6. 设 \(\mathfrak{g}\) 是李代数，E 是由 1 和 \(\mathcal{L}(\mathfrak{g})\) 生成的 \(\operatorname{ad}_{\mathfrak{g}}x(x\in \mathfrak{g})\) 的结合子代数。令 R 为 E 的 Jacobson 根。

\begin{enumerate}
\def\labelenumi{(\alph{enumi})}
\item
  满足 \(\mathfrak{n}\) 的 \(y\in \mathfrak{g}\) 的集合 \(\operatorname {ad}_{\mathfrak{g}}y\in \mathbb{R}\) 是 \(\mathfrak{g}\) 的最大幂零理想。
\item
  它关于 Killing 形式与 g 正交。
\end{enumerate}

应注意，g/n 可能有非零幂零理想。

